% Compile in Overleaf with XeLaTeX.
\documentclass[a4paper,12pt]{article}

\usepackage{fontspec}
\usepackage{polyglossia}
\setmainlanguage{russian}
\setotherlanguage{english}
\setmainfont{DejaVu Serif}[
  BoldFont={DejaVu Serif Bold},
  ItalicFont={DejaVu Serif},
  BoldItalicFont={DejaVu Serif Bold}
]
\setsansfont{DejaVu Sans}
\setmonofont{DejaVu Sans Mono}

\usepackage{amsmath,amssymb}
\usepackage{geometry}
\usepackage{graphicx}
\usepackage{booktabs}
\usepackage{tabularx}
\usepackage{array}
\usepackage{float}
\usepackage{caption}
\usepackage{subcaption}
\usepackage{enumitem}
\usepackage{xcolor}
\usepackage{listings}
\usepackage{microtype}
\usepackage{fancyhdr}
\usepackage[hidelinks]{hyperref}

\geometry{left=25mm,right=20mm,top=20mm,bottom=22mm}
\setlength{\parindent}{1.25cm}
\setlength{\parskip}{0.15em}
\emergencystretch=6em
\renewcommand{\arraystretch}{1.18}
\captionsetup{font=small,labelfont=bf}
\setlist[itemize]{leftmargin=1.4cm,itemsep=0.2em}
\setlist[enumerate]{leftmargin=1.4cm,itemsep=0.2em}

\definecolor{codeblue}{RGB}{31,78,121}
\definecolor{codegreen}{RGB}{55,86,35}
\definecolor{codegray}{RGB}{245,245,245}
\lstset{
  language=Python,
  basicstyle=\ttfamily\footnotesize,
  keywordstyle=\color{codeblue}\bfseries,
  commentstyle=\color{codegreen},
  stringstyle=\color{purple},
  backgroundcolor=\color{codegray},
  frame=single,
  rulecolor=\color{gray!40},
  numbers=left,
  numberstyle=\tiny\color{gray},
  breaklines=true,
  showstringspaces=false,
  columns=fullflexible,
  keepspaces=true
}

\pagestyle{fancy}
\fancyhf{}
\fancyhead[L]{Алгоритмы компьютерной графики}
\fancyhead[R]{Лабораторная работа № 1}
\fancyfoot[C]{\thepage}
\setlength{\headheight}{15pt}

% ---------- Edit these fields before submission ----------
\newcommand{\studentname}{Чжун Ц. Су Л.}
\newcommand{\studentgroup}{указать группу}
\newcommand{\teachername}{указать преподавателя}
\newcommand{\reportyear}{2026}

\begin{document}

\begin{titlepage}
  \thispagestyle{empty}
  \begin{center}
    \textbf{МИНИСТЕРСТВО НАУКИ И ВЫСШЕГО ОБРАЗОВАНИЯ РФ}\\[0.6em]
    \textbf{УНИВЕРСИТЕТ ИТМО}\\[0.6em]
    Факультет программной инженерии и компьютерной техники

    \vfill

    {\Large\bfseries ОТЧЁТ ПО ЛАБОРАТОРНОЙ РАБОТЕ № 1}\\[1.2em]
    {\large по дисциплине «Алгоритмы компьютерной графики»}\\[2em]
    {\Large\bfseries Расчёт освещённости на плоскости\\
    от точечного источника света}

    \vfill

    \begin{tabular}{@{}ll@{}}
      Выполнил: & \studentname \\
      Группа: & \studentgroup \\
      Преподаватель: & \teachername
    \end{tabular}

    \vfill
    Санкт-Петербург\\
    \reportyear
  \end{center}
\end{titlepage}

\tableofcontents
\newpage

\section{Цель и постановка задачи}

Цель работы состоит в освоении методов расчёта и визуализации распределения
энергетической освещённости на плоскости, создаваемой точечным источником с
ламбертовской диаграммой излучения.

В работе требуется:
\begin{itemize}
  \item рассчитать распределение освещённости в пределах заданной окружности;
  \item сформировать изображение заданного разрешения с нормировкой значений
        в диапазон от 0 до 255;
  \item обеспечить изменение параметров модели в пользовательском интерфейсе;
  \item записать сформированное изображение в файл и вывести его на экран;
  \item построить график сечения, проходящего через центр области;
  \item определить освещённость в пяти контрольных точках, а также минимальное,
        максимальное и среднее значения внутри окружности.
\end{itemize}

Расчётная плоскость совпадает с плоскостью $z=0$. Источник расположен в точке
$L(x_L,y_L,z_L)$, причём его оптическая ось направлена по отрицательному
направлению оси $Z$. Прямоугольная область изображения имеет физические размеры
$W\times H$, а область расчёта ограничена окружностью с центром
$C(x_C,y_C)$ и радиусом $R$.

\section{Теоретические сведения}

\subsection{Основные радиометрические величины}

Поток излучения $\Phi$ определяет количество энергии $Q$, переносимой в единицу
времени:
\begin{equation}
  \Phi=\frac{dQ}{dt}, \qquad [\Phi]=\text{Вт}.
\end{equation}

Сила излучения $I$ характеризует поток в заданном направлении, отнесённый к
элементу телесного угла $d\Omega$:
\begin{equation}
  I=\frac{d\Phi}{d\Omega}, \qquad [I]=\text{Вт/ср}.
\end{equation}

Энергетическая освещённость $E$ представляет собой плотность потока,
падающего на поверхность:
\begin{equation}
  E=\frac{d\Phi}{dS}, \qquad [E]=\text{Вт/м}^{2}.
\end{equation}

Для точечного источника освещённость площадки на расстоянии $d$ определяется
законом обратных квадратов и косинусным законом:
\begin{equation}
  E=\frac{I(\theta)\cos\alpha}{d^{2}},
  \label{eq:inverse-square}
\end{equation}
где $\alpha$ --- угол между нормалью к принимающей поверхности и направлением
на источник.

\subsection{Ламбертовская диаграмма излучения}

Для ламбертовского источника сила излучения зависит от угла $\theta$ между
лучом и оптической осью по закону
\begin{equation}
  I(\theta)=I_0\cos\theta,
  \label{eq:lambert}
\end{equation}
где $I_0$ --- сила излучения в направлении оптической оси ($\theta=0$).

Важно различать два косинусных множителя: $\cos\theta$ в
формуле~\eqref{eq:lambert} описывает диаграмму источника, а $\cos\alpha$ в
формуле~\eqref{eq:inverse-square} учитывает проекцию принимающей площадки.

\subsection{Расчётная формула}

Рассмотрим точку плоскости $P(x,y,0)$. Квадрат расстояния от источника до
этой точки равен
\begin{equation}
  d^2=(x-x_L)^2+(y-y_L)^2+z_L^2.
  \label{eq:distance}
\end{equation}

Так как оптическая ось источника перпендикулярна плоскости, углы излучения и
падения численно равны:
\begin{equation}
  \cos\theta=\cos\alpha=\frac{z_L}{d}.
\end{equation}

Подставляя это соотношение в формулы~\eqref{eq:inverse-square} и
\eqref{eq:lambert}, получим итоговую формулу:
\begin{equation}
  \boxed{
  E(x,y)=
  \frac{I_0z_L^2}
  {\left((x-x_L)^2+(y-y_L)^2+z_L^2\right)^2}}
  \label{eq:result}
\end{equation}

В формуле~\eqref{eq:result} расстояния должны быть выражены в метрах. Если
координаты непосредственно подставляются в миллиметрах, результат переводится
в Вт/м$^2$ умножением на $10^6$:
\begin{equation}
  E_{\text{Вт/м}^2}=10^6
  \frac{I_0z_{L,\text{мм}}^2}{d_{\text{мм}}^4}.
\end{equation}

Поскольку исходная сила излучения задана в Вт/ср, в работе рассчитывается
энергетическая освещённость в Вт/м$^2$, а не фотометрическая освещённость в
люксах. Для перехода к фотометрическим величинам необходимо знать спектральное
распределение источника.

\section{Исходные данные}

Для демонстрационного варианта выбраны параметры, приведённые в
таблице~\ref{tab:parameters}.

\begin{table}[H]
  \centering
  \caption{Параметры расчётного варианта}
  \label{tab:parameters}
  \begin{tabular}{@{}lll@{}}
    \toprule
    Параметр & Обозначение & Значение \\
    \midrule
    Ширина области & $W$ & 1000 мм \\
    Высота области & $H$ & 1000 мм \\
    Разрешение & $W_{res}\times H_{res}$ & $600\times600$ пикселей \\
    Координаты источника & $(x_L,y_L,z_L)$ & $(200,-150,800)$ мм \\
    Осевая сила излучения & $I_0$ & 100 Вт/ср \\
    Центр окружности & $(x_C,y_C)$ & $(0,0)$ мм \\
    Радиус окружности & $R$ & 400 мм \\
    \bottomrule
  \end{tabular}
\end{table}

Размер одного пикселя по обеим осям равен
\begin{equation}
  \Delta x=\frac{W}{W_{res}}=1{,}6667\text{ мм},\qquad
  \Delta y=\frac{H}{H_{res}}=1{,}6667\text{ мм}.
\end{equation}
Следовательно, условие квадратных пикселей $\Delta x=\Delta y$ выполнено.

\section{Алгоритм программы}

Алгоритм расчёта включает следующие этапы:
\begin{enumerate}
  \item Ввод физических размеров изображения, разрешения, положения
        источника, силы излучения, центра и радиуса окружности.
  \item Проверка допустимых диапазонов и условия квадратных пикселей
        $W/W_{res}=H/H_{res}$.
  \item Формирование массивов координат центров пикселей:
  \begin{equation}
    x_j=-\frac{W}{2}+\left(j+\frac12\right)\frac{W}{W_{res}},
    \qquad
    y_i=-\frac{H}{2}+\left(i+\frac12\right)\frac{H}{H_{res}}.
  \end{equation}
  \item Формирование логической маски окружности:
  \begin{equation}
    (x-x_C)^2+(y-y_C)^2\le R^2.
  \end{equation}
  \item Вычисление $E(x,y)$ по формуле~\eqref{eq:result} для всех центров
        пикселей с использованием векторных операций NumPy.
  \item Вычисление $E_{\min}$, $E_{\max}$ и $E_{\text{ср}}$ только по
        пикселям, находящимся внутри окружности.
  \item Нормировка значений внутри окружности:
  \begin{equation}
    G(x,y)=\operatorname{round}\left(255\frac{E(x,y)}{E_{\max}}\right).
  \end{equation}
        За пределами окружности принимается $G=0$.
  \item Сохранение восьмибитного изображения, цветной карты распределения,
        графика сечения, таблицы пяти точек и файла с численными результатами.
\end{enumerate}

Вычислительная сложность алгоритма составляет
$O(W_{res}H_{res})$, а объём используемой памяти также пропорционален числу
пикселей. При максимальном разрешении $800\times800$ вычисляется 640000 точек,
что позволяет выполнять расчёт интерактивно.

Ключевая часть программной реализации имеет вид:
\begin{lstlisting}[caption={Вычисление распределения и нормировка}]
dx_m = (grid_x_mm - source_x_mm) / 1000.0
dy_m = (grid_y_mm - source_y_mm) / 1000.0
z_m = source_z_mm / 1000.0
d2_m2 = dx_m**2 + dy_m**2 + z_m**2
E = I0 * z_m**2 / d2_m2**2

mask = (grid_x_mm - circle_x_mm)**2 \
     + (grid_y_mm - circle_y_mm)**2 <= radius_mm**2
normalized[mask] = np.rint(255 * E[mask] / E[mask].max())
\end{lstlisting}

Полный исходный код находится в приложенном файле
\path{lab1_illumination.py}. Программа поддерживает графический режим и
режим воспроизводимого запуска без интерфейса:
\begin{center}
  \texttt{python lab1\_illumination.py --no-gui --output results}
\end{center}

\section{Результаты}

\subsection{Распределение освещённости}

На рисунке~\ref{fig:map} показано рассчитанное распределение. Максимум
находится около проекции источника $(200,-150)$ мм, которая расположена внутри
расчётной окружности. По мере увеличения горизонтального расстояния до этой
точки освещённость уменьшается.

\begin{figure}[H]
  \centering
  \includegraphics[width=0.82\textwidth]{results/illumination_map.png}
  \caption{Распределение энергетической освещённости внутри окружности}
  \label{fig:map}
\end{figure}

Сформированное восьмибитное изображение сохранено в файле
\path{results/illumination_normalized.png}. Значение 255 соответствует
максимальной освещённости выбранного варианта, а нулевые пиксели за пределами
окружности служат фоном.

\subsection{Сечение через центр области}

Сечение рассчитано непосредственно по физической формуле вдоль прямой
$y=y_C=0$ мм при $x\in[-400,400]$ мм. Это исключает погрешность, которая могла
бы возникнуть при выборе ближайшей строки дискретного изображения.

\begin{figure}[H]
  \centering
  \includegraphics[width=0.86\textwidth]{results/center_section.png}
  \caption{Горизонтальное сечение распределения через центр окружности}
  \label{fig:section}
\end{figure}

Максимум кривой расположен при $x=200$ мм, поскольку в этой точке проходит
проекция источника на выбранную линию сечения.

\subsection{Контрольные значения}

Под пятью контрольными точками понимаются центр окружности и четыре точки её
границы в положительных и отрицательных направлениях осей $X$ и $Y$. Значения
вычислены непосредственно по формуле~\eqref{eq:result}, а не выбраны из
ближайших пикселей.

\begin{table}[H]
  \centering
  \caption{Освещённость в пяти контрольных точках}
  \label{tab:points}
  \begin{tabular}{@{}p{38mm}r@{\hspace{7mm}}r@{\hspace{7mm}}r@{}}
    \toprule
    Точка & $x$, мм & $y$, мм & $E$, Вт/м$^2$ \\
    \midrule
    Центр $C$ & 0 & 0 & 129{,}684275 \\
    Граница $+X$ & 400 & 0 & 129{,}684275 \\
    Граница $-X$ & -400 & 0 & 61{,}214364 \\
    Граница $+Y$ & 0 & 400 & 66{,}300203 \\
    Граница $-Y$ & 0 & -400 & 116{,}087928 \\
    \bottomrule
  \end{tabular}
\end{table}

Статистические значения по центрам пикселей внутри окружности представлены в
таблице~\ref{tab:statistics}.

\begin{table}[H]
  \centering
  \caption{Статистика освещённости внутри области расчёта}
  \label{tab:statistics}
  \begin{tabular}{@{}lr@{}}
    \toprule
    Характеристика & Значение, Вт/м$^2$ \\
    \midrule
    Минимальное значение $E_{\min}$ & 56{,}705531 \\
    Максимальное значение $E_{\max}$ & 156{,}249322 \\
    Среднее значение $E_{\text{ср}}$ & 110{,}434019 \\
    \bottomrule
  \end{tabular}
\end{table}

Непрерывные аналитические границы для того же варианта равны
$56{,}692042$ и $156{,}250000$ Вт/м$^2$. Небольшое отличие дискретных
экстремумов объясняется тем, что координаты экстремальных точек не обязаны
совпадать с центрами пикселей.

\section{Проверка корректности}

Для проверки программы использованы следующие свойства модели:
\begin{enumerate}
  \item При расположении источника над центром окружности распределение должно
        быть осесимметричным, а максимум должен находиться в центре.
  \item Непосредственно под источником $d=z_L$, поэтому
        $E=I_0/z_L^2$.
  \item При умножении $I_0$ на коэффициент $k$ все физические значения $E$
        умножаются на $k$, но изображение после нормировки на собственный
        максимум не изменяется.
  \item При увеличении разрешения геометрическая форма распределения
        сохраняется, а дискретные экстремумы приближаются к аналитическим.
  \item Для выбранного варианта аналитический максимум равен
        $100/0{,}8^2=156{,}25$ Вт/м$^2$, что согласуется с результатом программы.
\end{enumerate}

\section{Вывод}

В ходе работы реализован расчёт энергетической освещённости горизонтальной
плоскости от точечного источника с ламбертовской диаграммой. В программе
учтены закон обратных квадратов, косинусная диаграмма источника и косинус
угла падения на плоскость. Реализованы проверка параметров, построение
физической координатной сетки с квадратными пикселями, ограничение области
окружностью, нормировка в диапазон 0--255, сохранение изображения, вычисление
контрольных значений и построение сечения через центр области.

Полученные численные результаты согласуются с аналитическими оценками.
Расхождение между аналитическими и дискретными экстремумами определяется
конечным пространственным разрешением изображения.

\begin{thebibliography}{9}
  \bibitem{zhdanov}
  Жданов Д.~Д., Потемин И.~С., Жданов А.~Д.
  \textit{Теоретические основы компьютерной графики и вычислительной оптики}.
  СПб.: Университет ИТМО, 2023.

  \bibitem{ignatenko}
  Игнатенко А.
  \textit{Свет как энергия. Радиометрия. Фотометрия}: материалы лекции.
  МГУ, 2011.

  \bibitem{assignment}
  Лабораторные работы по курсу «Алгоритмы компьютерной графики».
  ЛР~1 «Расчёт освещённости на плоскости от точечного источника света».
\end{thebibliography}

\end{document}
